\documentclass{article}
\usepackage[T1]{fontenc}
\usepackage{lmodern}
\usepackage{graphicx}
\usepackage{amsmath,amssymb}
\usepackage[a4paper,margin=2cm]{geometry}
\usepackage[absolute,overlay]{textpos}
\setlength{\TPHorizModule}{1mm}
\setlength{\TPVertModule}{1mm}
\usepackage[indonesian]{babel}
\usepackage{hyperref}
\setlength{\emergencystretch}{2em}
% unit: Halaman 3

\begin{document}

% === Halaman 3 (llm) ===

\section*{Hybrid Cryptography}

\begin{itemize}
    \item Alice dan Bob sepakat mengenkripsi dan mendekripsi pesan dengan algoritma kriptografi simetri (misalkan AES).
    \item Bagaimana cara Bob dapat mengetahui kunci AES yang digunakan oleh Alice (kunci AES)?
    \item Solusinya adalah dengan menggunakan \textit{hybrid cryptography}.
    \item \textit{Hybrid cryptography}: menggabungkan kriptografi kunci-simetri (misalkan AES) dengan kriptografi kunci-publik (misalkan RSA).
    \item Caranya sebagai berikut:
    \begin{enumerate}
        \item Alice memiliki sepasang kunci privat ($\text{SK}_{\text{Alice}}$) dan kunci publik ($\text{PK}_{\text{Alice}}$) miliknya.
        \item Bob juga memiliki sepasang kunci privat ($\text{SK}_{\text{Bob}}$) dan kunci publik ($\text{PK}_{\text{Bob}}$) miliknya.
    \end{enumerate}
\end{itemize}

\end{document}
